\textbf{COLUNI 2024 -- Questão 07}\par\medskip Observe o seguinte trecho da \textit{Carta pras Icamiabas}, de \textit{Macunaíma}.

\par

“Andam \underline{elas} vestidas de rutilantes jóias e panos finíssimos, que lhes acentuam o donaire do porte, e mal encobrem as graças, que, a de nenhuma outra cedem pelo formoso do torneado e pelo tom. São sempre alvíssimas \underline{as donas de cá}; e tais e tantas habilidades demonstram, no brincar, que enumerá-las, aqui seria fastiendo porventura; e, certamente, quebraria os mandamentos de discreção, que em relação de Imperator para súbditas se requer. Que beldades! Que elegância! Que cachet! Que degagé flamífero, ignívomo, devorador!! Só pensamos \underline{nelas}, muito embora não nos descuidemos, relapso, da \underline{nossa muiraquitã}”. (ANDRADE, Mário. 2001, p.73)

\par

A partir do contexto da carta, e da análise do excerto, os trechos destacados referem-se, respectivamente, a:\begin{enumerate}
\item paulistanas, paulistanas, icamiabas, paulistanas.
\item paulistanas, icamiabas, icamiabas, paulistanas.
\item paulistanas, paulistanas, paulistanas, paulistanas.
\item icamiabas, icamiabas, paulistanas, paulistanas.\end{enumerate}
